\documentclass[../../main.tex]{subfiles}

\begin{document}

\section{The Multiplicative Inverse}

As we discuss matrix multiplication, you may have wondered about
matrix division. I mean it would be so nice if we could just
divide the coefficient matrix to find our variables. Although there
is no ``division'' operation per se, we could achieve similar
operation with the multiplicative inverse of a matrix, or simply
the inverse matrix. In other words, some solutions to linear systems
can be written as $x = A^{-1}b$.

\begin{definition}{}{}
The \textbf{inverse of a matrix} is a matrix that produces an identity
matrix when multiplied with the original matrix. The inverse of
$A$ is denoted as $A^{-1}$.
\end{definition}

We could see that matrices $A$ and $B$ are the inverse of each other
if the following equation is satisfied.
\[
AB = BA = I
\]
Below is an example where $B = A^{-1}$.
\[
A =
\begin{bmatrix}
    1 & 2 \\
    3 & 5
\end{bmatrix},
\quad
B =
\begin{bmatrix}
    -5 & 2 \\
    3  & -1
\end{bmatrix}
\]
One thing to note from the definition is that $A$ and $B$ must be
square matrices. Say we have $A_{4 \times 2}$ and $B_{2 \times 4}$.
Even if both $AB$ and $BA$ return an identity matrix, the order will
be different. Therefore, for a matrix to have the inverse, it must
be a square matrix.

\begin{definition}{}{}
A matrix is \textbf{invertible}, or \textbf{nonsingular}, if it
retains an inverse matrix. A matrix is \textbf{singular} if it
does not have an inverse matrix.
\end{definition}

Now from the definition above, we could establish a theorem.

\begin{theorem}{}{Diagonal Elements and Invertibility}
If all diagonal elements of a matrix are nonzero constants,
then the matrix is invertible.
\end{theorem}
\begin{proof}
Consider the matrix $A_{m \times m}$ defined below.
\[
A =
\begin{bmatrix}
    a_{11} & a_{12} & a_{13} & \cdots & a_{1m} \\
    a_{21} & a_{22} & a_{23} & \cdots & a_{2m} \\
    a_{31} & a_{32} & a_{33} & \cdots & a_{3m} \\
    \vdots & \vdots & \vdots & \ddots & \vdots \\
    a_{m1} & a_{m2} & a_{m3} & \cdots & a_{mm}
\end{bmatrix}
\]
It suffices to show that there exists the inverse matrix
for all values of the elements. Consider the inverse below.
\[
A^{-1} =
\begin{bmatrix}
    \frac{1}{a_{11}} & 0 & 0 & \cdots & 0 \\
    0 & \frac{1}{a_{22}} & 0 & \cdots & 0 \\
    0 & 0 & \frac{1}{a_{33}} & \cdots & 0 \\
    \vdots & \vdots & \vdots & \ddots & \vdots \\
    0 & 0 & 0 & \cdots & \frac{1}{a_{mm}}
\end{bmatrix}
\]
Because $a_{ii}$ for integer $i \in [1, m]$ are nonzero numbers,
all elements in the matrix $A^{-1}$ are defined, and the inverse
matrix exists.
\end{proof}

Now, it is important to note that not all square matrices have the
inverse. For example, consider a matrix with zero row. The same
row in the product matrices will become a zero row. By definition,
it is not possible to make an identity matrix with a zero row.

Now that we discussed the existence, we could also investigate if
the inverse matrix is unique just as a nonzero integer has a unique
reciprocal.

\begin{theorem}{}{Uniqueness of Inverse Matrix}
If there exists an inverse matrix, the inverse is unique.
\end{theorem}
\begin{proof}
For the sake of contradiction, let $B$ and $C$ be different
inverse of $A$. Consider the following equation.
\[
B = BI = B(AC) = (BA)C = IC = C
\]
Because the fact that $B = C$ contradicts the initial assumption
that $B \neq C$, all inverse matrices are unique.
\end{proof}

Before we discuss how we actually find an inverse, let's discuss
four important properties of inverse matrices.

\begin{theorem}{}{Properties of the Inverse Matrix}
For invertible square matrices $A$ and $B$ and a nonzero scalar
$\lambda$, the following equations hold.
\begin{enumerate}
\item $(A^{-1})^{-1} = A$
\item $(A^\intercal)^{-1} = (A^{-1})^\intercal$
\item $(AB)^{-1} = B^{-1}A^{-1}$
\item $(\lambda A)^{-1} = (\frac{1}{\lambda})A^{-1}$
\end{enumerate}
\end{theorem}

There's a lot to prove and let's do them one by one starting with
the first one!

\begin{proof}
Let matrix $B = A^{-1}$. Substituting, it suffices to show that
$B^{-1} = A$. By definition of inverse matrix, the equation
$AB = BA = 1$ holds. In other words, $A$ is the inverse of $B$
and the property holds.
\end{proof}

Below is the proof for the second property.

\begin{proof}
Notice that $(AB)^\intercal = B^\intercal A^\intercal$ from
Theorem \ref{theo:Properties of Transpose of a Matrix}. Let matrix
$B = A^{-1}$ and consider the following equation.
\[
I
= (AB)^\intercal
= B^\intercal A^\intercal
= \left( A^{-1} \right)^\intercal A^\intercal
\]
In other words, $(A^{-1})^\intercal$ is the inverse of $A^\intercal$,
or $(A^\intercal)^{-1}$, proving the property.
\end{proof}

Continuing, here is the proof for the third property.

\begin{proof}
Consider the following equation that uses the associative property
of matrix multiplication.
\[
(B^{-1}A^{-1})(AB)
= B^{-1}(A^{-1}A)B
= B^{-1}(I)B
= B^{-1}(IB)
= B^{-1}B
= I
\]
Because $B^{-1}A^{-1}$ is the inverse of $AB$, the property is true.
\end{proof}

Finally, here is the proof for the last property.

\begin{proof}
Let $B = \lambda A$. Consider the following equation.
\[
\frac{1}{\lambda} A^{-1} B
= \frac{1}{\lambda} A^{-1} \lambda A
= A^{-1} A
= I
\]
In other words, $\frac{1}{\lambda} A^{-1}$ is the inverse of
$\lambda A$, proving the fourth property.
\end{proof}

Now that we took a look at different properties of inverse matrices,
let's discuss how to actually find one with tools to consider.

\begin{definition}{}{}
For a matrix $A$ that is not necessarily a square matrix, an
\textbf{elementary matrix}, denoted as $E$, is a square matrix that
performs an elementary row operation by multiplying with $A$.
\end{definition}

Let's take a look at an example. Consider matrix $A$ defined below,
and we want to perform elementary row operations of changing the
order of rows $1$ and $3$ and add two times the fourth row to the
second respectively.
\[
A =
\begin{bmatrix}
    1 & 2 & 3 & 4 \\
    1 & 0 & 2 & 0 \\
    2 & 4 & 6 & 8 \\
    1 & 1 & 1 & 1
\end{bmatrix}
\]
The elementary row operations above can be complicated with words,
but we can represent that with two elementary matrices respectively.
Letting $E_1$ and $E_2$ be elementary matrices for each operation,
we can represent the elementary row operations as the following.
\[
E_1 =
\begin{bmatrix}
    0 & 0 & 1 & 0 \\
    0 & 1 & 0 & 0 \\
    1 & 0 & 0 & 0 \\
    0 & 0 & 0 & 1
\end{bmatrix},
\quad
E_2 =
\begin{bmatrix}
    1 & 0 & 0 & 0 \\
    0 & 1 & 0 & 2 \\
    0 & 0 & 1 & 0 \\
    0 & 0 & 0 & 1
\end{bmatrix}
\]
Let's multiply to see if we are right.
\[
E_1 A =
\begin{bmatrix}
    \textcolor{Red}{2} &
    \textcolor{Red}{4} &
    \textcolor{Red}{6} &
    \textcolor{Red}{8} \\
    1 & 0 & 2 & 0 \\
    \textcolor{Red}{1} &
    \textcolor{Red}{2} &
    \textcolor{Red}{3} &
    \textcolor{Red}{4} \\
    1 & 1 & 1 & 1
\end{bmatrix},
\quad
E_2 A =
\begin{bmatrix}
    1 & 2 & 3 & 4 \\
    \textcolor{Red}{3} &
    \textcolor{Red}{2} &
    \textcolor{Red}{4} &
    \textcolor{Red}{2} \\
    2 & 4 & 6 & 8 \\
    1 & 1 & 1 & 1
\end{bmatrix}
\]
Now we can notice three observations and tips for constructing
elementary matrices. For each elementary row operation, we could
do the following.

\begin{definition}{}{}
\textbf{Methods to construct the elementary matrix for each
elementary row operation:}
\begin{itemize}[noitemsep]
\item
To interchange row $a$ and $b$, write the corresponding identity
matrix and interchange row $a$ and $b$ from the identity matrix.
\item
To multiply a nonzero scalar $k$ to $a^\text{th}$ row, start from the
corresponding identity matrix and change the $1$ from $a^\text{th}$
row of the identity matrix to $k$.
\item
To add row $a$ multiplied by the nonzero scalar $k$ to row $b$,
start from the corresponding identity matrix and find the
$b^\text{th}$ row in the matrix. Then, change the zero from the
$a^\text{th}$ column of the row to $k$.
\end{itemize}
\end{definition}

From this, we can notice the following theorem.

\begin{theorem}{}{Invertibility of Elementary Matrices}
All elementary matrices are invertible.
\end{theorem}
\begin{proof}
The proof is rather simple. Because all elementary row operations
can be inverted with the inverse operation, all elementary matrices
are invertible.
\end{proof}

Building onto the theorem, we could prove the uniqueness.

\begin{theorem}{}{Uniqueness of Elementary Matrices}
For each elementary row operation, the corresponding elementary
matrix is unique.
\end{theorem}
\begin{proof}
For the sake of contradiction, assume there exist elementary
matrices $E$ and $F \neq E$. By definition, $EA = FA$ and
$A = EE^{-1}A = E^{-1}FA$. By Theorem \ref{theo:Uniqueness of Identity
Matrices}, $E^{-1}F = I$ and $E = F$ is obtained by Theorem
\ref{theo:Uniqueness of Inverse Matrix}. By contradiction, it is
shown that elementary matrices are unique.
\end{proof}

There are more interesting lemmas and theorems, but let's discuss
them in latter notes. However, let's discuss a few theorems for each
elementary row operation before moving on to finding an inverse
matrix.

\begin{theorem}{}{Three Properties of Elementary Matrices}
For each elementary matrix for corresponding elementary row
operations, the following properties hold.
\begin{enumerate}[noitemsep]
\item
For an elementary matrix that changes the order of two rows,
the inverse of such a matrix is itself.
\item
For an elementary matrix that scales a row with nonzero constant
$k$, the inverse can be obtained by replacing $k$ in the elementary
matrix with $\frac{1}{k}$.
\item
For an elementary matrix that adds the scaled row into the other,
its inverse can be obtained by replacing the nonzero scalar $k$
with $-k$.
\end{enumerate}
\end{theorem}

As always, let's start by proving the first property.

\begin{proof}
The inverse of an elementary matrix can be considered as undoing
the elementary row operation. To undo the interchanging of two rows,
the two rows can be interchanged again. In other words, because
the elementary row operation $EE = I$ by Theorem \ref{theo:Uniqueness
of Identity Matrices}, $E = E^{-1}$.
\end{proof}

Continuing, here is the proof for the second property.

\begin{proof}
Let elementary matrices $E$ and $F$ be matrices that scale a
row by a nonzero scalar $k$ and $\frac{1}{k}$ respectively.
By definition, $F(EI) = I$ and $FE = I$. Therefore, $F = E^{-1}$
and the property is shown.
\end{proof}

Lastly, here is the proof for the final property.

\begin{proof}
Similar to the proof for the second property, let elementary matrices
$E$ and $F$ represent the elementary row operations that add a row
scaled by nonzero constant $k$ to another and the operation that adds
a row scaled by $-k$ to another respectively. By definition,
$F(EI) = I$ and $F = E^{-1}$ as demonstrated above.
\end{proof}

Now before we actually find an inverse, we need to know which
square matrices have the inverse matrix.

\begin{definition}{}{}
A square matrix $A$ is invertible if and only if it can be represented
as an identity matrix after sequences of elementary row operations,
i.e. there exist elementary matrices $E_i$ such that the following
equation is true.
\[
E_k E_{k-1} \cdots E_1 A = I
\]
\end{definition}

This is a very important result that we will prove later in a note,
so don't worry. Finally, here are the steps to find an inverse if
there exists one.

\begin{definition}{}{}
\textbf{Steps to Find the Inverse Matrix of
$\mathbf{A_{m \times m}}$:}
\begin{enumerate}[noitemsep]
\item
Write an augmented matrix $[A \mid I]$ for $I_{m}$.
\item
Utilize elementary row operations on the augmented matrix for
the block $A$ to be in REF.
\item
If the main diagonal of the left partition contains a zero,
$A$ is singular. If not, the next step can be continued.
\item
Use elementary row operations again to transform the block $A$
in REF to an identity matrix. The right partition block is
the inverse of $A$.
\end{enumerate}
\end{definition}

Let's take a look at a few examples.

\begin{problem}{}{}
Find the inverse of $A =
\begin{bmatrix}
    1 & 2 & 3 \\
    4 & 5 & 6 \\
    7 & 8 & 9
\end{bmatrix}$.
\end{problem}
\begin{solution}
First, the augmented matrix can be constructed.
\[
\begin{bNiceArray}{ccc|ccc}
    1 & 2 & 3 & 1 & 0 & 0 \\
    4 & 5 & 6 & 0 & 1 & 0 \\
    7 & 8 & 9 & 0 & 0 & 1
\end{bNiceArray}
\]
Consider the following elementary row operations.
\begin{align*}
\begin{bNiceArray}{ccc|ccc}
    1 & 2 & 3 & 1 & 0 & 0 \\
    4 & 5 & 6 & 0 & 1 & 0 \\
    7 & 8 & 9 & 0 & 0 & 1
\end{bNiceArray}
&\rightarrow
\begin{bNiceArray}{ccc|ccc}
    1 & 2 & 3 & 1 & 0 & 0 \\
    8 & 10 & 12 & 0 & 2 & 0 \\
    7 & 8 & 9 & 0 & 0 & 1
\end{bNiceArray}
\displaybreak
\rightarrow
\begin{bNiceArray}{ccc|ccc}
    1 & 2 & 3 & 1 & 0 & 0 \\
    1 & 2 & 3 & 0 & 2 & -1 \\
    7 & 8 & 9 & 0 & 0 & 1
\end{bNiceArray} \\
&\rightarrow
\begin{bNiceArray}{ccc|ccc}
    1 & 2 & 3 & 1 & 0 & 0 \\
    0 & 0 & 0 & -1 & 2 & -1 \\
    7 & 8 & 9 & 0 & 0 & 1
\end{bNiceArray}
\end{align*}
It is evident that the main diagonal will contain a zero after
the transformation. Therefore, the matrix is not invertible.
\end{solution}

\begin{problem}{}{}
Find the inverse of the following matrix $A$.
\[
A =
\begin{bmatrix}
    1 & 0 & 0 \\
    2 & 1 & 1 \\
    0 & 1 & 2
\end{bmatrix}
\]
\end{problem}
\begin{solution}
First and foremost, the augmented matrix can be constructed.
\begin{align*}
\begin{bNiceArray}{ccc|ccc}
    1 & 0 & 0 & 1 & 0 & 0 \\
    2 & 1 & 1 & 0 & 1 & 0 \\
    0 & 1 & 2 & 0 & 0 & 1
\end{bNiceArray}
&\rightarrow
\begin{bNiceArray}{ccc|ccc}
    1 & 0 & 0 & 1 & 0 & 0 \\
    0 & 1 & 2 & 0 & 0 & 1 \\
    2 & 1 & 1 & 0 & 1 & 0
\end{bNiceArray}
\rightarrow
\begin{bNiceArray}{ccc|ccc}
    1 & 0 & 0 & 1 & 0 & 0 \\
    0 & 1 & 2 & 0 & 0 & 1 \\
    0 & 0 & -1 & -2 & 1 & -1
\end{bNiceArray} \\
&\rightarrow
\begin{bNiceArray}{ccc|ccc}
    1 & 0 & 0 & 1 & 0 & 0 \\
    0 & 1 & 2 & 0 & 0 & 1 \\
    0 & 0 & 1 & 2 & -1 & 1
\end{bNiceArray}
\end{align*}
Continuing, the left partition can be transformed into an
identity matrix with elementary row operations.
\[
\begin{bNiceArray}{ccc|ccc}
    1 & 0 & 0 & 1 & 0 & 0 \\
    0 & 1 & 2 & 0 & 0 & 1 \\
    0 & 0 & 1 & 2 & -1 & 1
\end{bNiceArray}
\rightarrow
\begin{bNiceArray}{ccc|ccc}
    1 & 0 & 0 & 1 & 0 & 0 \\
    0 & 1 & 0 & -4 & 2 & -1 \\
    0 & 0 & 1 & 2 & -1 & 1
\end{bNiceArray}
\]
Therefore, the inverse of $A$ is the following.
\[
A^{-1} =
\begin{bmatrix}
    1  & 0  & 0 \\
    -4 & 2  & -1 \\
    2  & -1 & 1
\end{bmatrix}
\]
\end{solution}

We could actually multiply to see if we get $I_3$.
\[
\begin{bmatrix}
    1 & 0 & 0 \\
    2 & 1 & 1 \\
    0 & 1 & 2
\end{bmatrix}
\begin{bmatrix}
    1  & 0  & 0 \\
    -4 & 2  & -1 \\
    2  & -1 & 1
\end{bmatrix}
=
\begin{bmatrix}
    1 & 0 & 0 \\
    0 & 1 & 0 \\
    0 & 0 & 1
\end{bmatrix}
\]
Indeed it is the inverse! In the next section of the note,
we will discuss LU decomposition before completing our note on
linear systems.

\end{document}


